% Auto-generated from run_prompt_experiments.py
% v5 here maps to script version b.v6 as requested.
\section{Appendix A: Prompt Templates (Scenario B)}
\label{app:prompt_templates_b}

This appendix reports prompt templates in boxed tables. We list versions v0--v5, where v5 corresponds to script version \texttt{b.v6}.

\begin{table*}[t]
  \centering
  \caption{Shared System Prompt for Scenario B prompt versions.}
  \label{tab:app_prompt_system}
  \fbox{\begin{minipage}{0.97\textwidth}
\begin{verbatim}
你是理性经济主体，目标是在不确定他人行为的情况下最大化你的期望效用。

【重要】你必须只输出一个有效的JSON对象，不要包含任何额外的文本、解释或markdown标记。
JSON必须包含 "share" 和 "reason" 两个字段。
确保 "reason" 字段的字符串正确闭合（以引号结束）。
\end{verbatim}
  \end{minipage}}
\end{table*}

\begin{table*}[t]
  \centering
  \caption{Prompt template v0 (source: \texttt{b.v0}).}
  \label{tab:app_prompt_v0}
  \fbox{\begin{minipage}{0.97\textwidth}
\begin{verbatim}
# 数据市场决策

你是用户 {user_id}，正在参与一个数据市场。

## 你的信息

- **报价**：p[{user_id}] = {price:.4f}
- **隐私偏好**：v[{user_id}] = {v_i:.3f}（单位信息成本）

## 决策

如果分享数据：
- 获得补偿 p = {price:.4f}
- 产生隐私成本 = v × 边际泄露量

**请权衡：补偿 vs 隐私成本**

## 输出

请输出严格JSON：
{{
  "share": 0或1（0=不分享，1=分享），
  "reason": "简要说明决策理由（不超过150字）"
}}
\end{verbatim}
  \end{minipage}}
\end{table*}

\begin{table*}[t]
  \centering
  \caption{Prompt template v1 (source: \texttt{b.v1}).}
  \label{tab:app_prompt_v1}
  \fbox{\begin{minipage}{0.97\textwidth}
\begin{verbatim}
# 数据市场决策

你是用户 {user_id}，正在参与一个数据市场。

## 你的私有信息

- **报价**：p[{user_id}] = {price:.4f}
- **隐私偏好**：v[{user_id}] = {v_i:.3f}（单位信息成本）

## 市场环境

- 用户总数：n = {n}
- 信息相关系数：ρ = {rho:.2f}
- 观测噪声：σ² = {sigma_noise_sq}
- 隐私偏好分布：所有用户的v ∈ [{v_min}, {v_max}]

你的位置：v = {v_i:.3f}（{v_description}）

## 决策框架

如果分享数据：
- 获得补偿 p = {price:.4f}
- 产生隐私成本 = v × 边际泄露量

**请权衡：补偿 vs 隐私成本**

## 输出

请输出严格JSON：
{{
  "share": 0或1（0=不分享，1=分享），
  "reason": "简要说明决策理由（不超过150字）"
}}
\end{verbatim}
  \end{minipage}}
\end{table*}

\begin{table*}[t]
  \centering
  \caption{Prompt template v2 (source: \texttt{b.v2}).}
  \label{tab:app_prompt_v2}
  \fbox{\begin{minipage}{0.97\textwidth}
\begin{verbatim}
# 数据市场决策

你是用户 {user_id}，正在参与一个数据市场。

## 你的私有信息

- **报价**：p[{user_id}] = {price:.4f}
- **隐私偏好**：v[{user_id}] = {v_i:.3f}（单位信息成本）

## 市场环境

**用户总数**：n = {n}

**信息相关系数**：ρ = {rho:.2f}
- 你的类型与其他用户的类型相关
- ρ = 0：他人信息无法推断你
- ρ = 1：他人信息完美推断你
- 当前 ρ = {rho:.2f}，推断能力{rho_level}

**观测噪声**：σ² = {sigma_noise_sq}
- σ² 越大：数据噪声越大，泄露程度越低
- σ² 越小：数据越准确，泄露程度越高

**隐私偏好分布**：v ∈ [{v_min}, {v_max}]
- 你的位置：v = {v_i:.3f}（{v_description}）

## 决策框架

如果分享数据：
- 获得补偿 p = {price:.4f}
- 产生隐私成本 = v × 边际泄露量

**请权衡：补偿 vs 隐私成本**

## 输出

请输出严格JSON：
{{
  "share": 0或1（0=不分享，1=分享），
  "reason": "简要说明决策理由（不超过150字）"
}}
\end{verbatim}
  \end{minipage}}
\end{table*}

\begin{table*}[t]
  \centering
  \caption{Prompt template v3 (source: \texttt{b.v3}).}
  \label{tab:app_prompt_v3}
  \fbox{\begin{minipage}{0.97\textwidth}
\begin{verbatim}
# 数据市场决策（推断外部性）

你是用户 {user_id}，正在参与一个数据市场。

## 你的私有信息

- **报价**：p[{user_id}] = {price:.4f}
- **隐私偏好**：v[{user_id}] = {v_i:.3f}（单位信息成本）

## 市场环境

**用户总数**：n = {n}

**信息相关系数**：ρ = {rho:.2f}
- 你的类型与其他用户的类型相关
- ρ = 0：他人信息无法推断你
- ρ = 1：他人信息完美推断你
- 当前 ρ = {rho:.2f}，推断能力{rho_level}

**观测噪声**：σ² = {sigma_noise_sq}
- σ² 越大：数据噪声越大，泄露程度越低
- σ² 越小：数据越准确，泄露程度越高

**隐私偏好分布**：v ∈ [{v_min}, {v_max}]
- 你的位置：v = {v_i:.3f}（{v_description}）

## 关键机制

**推断泄露**：即使你不分享，平台也能通过其他人的数据推断你的信息

**两种泄露**：
- **基础泄露**：其他人分享导致你的信息间接泄露
- **边际泄露**：你自己分享带来的额外泄露

**分享的真正成本 = 边际泄露**

## 决策框架

**如果分享**：
- 获得补偿：p = {price:.4f}
- 隐私成本：v × 边际泄露量
- 你的信息从部分泄露变为完全泄露

**如果不分享**：
- 无补偿
- 保护未间接泄露的部分信息
- 但仍有基础泄露

**请权衡：补偿 vs 边际隐私成本**

## 输出

请输出严格JSON：
{{
  "share": 0或1（0=不分享，1=分享），
  "reason": "简要说明决策理由（不超过150字）"
}}
\end{verbatim}
  \end{minipage}}
\end{table*}

\begin{table*}[t]
  \centering
  \caption{Prompt template v4 (source: \texttt{b.v4}).}
  \label{tab:app_prompt_v4}
  \fbox{\begin{minipage}{0.97\textwidth}
\begin{verbatim}
# 数据市场决策（推断外部性）

你是用户 {user_id}，正在参与一个数据市场。

## 你的私有信息

- **报价**：p[{user_id}] = {price:.4f}
- **隐私偏好**：v[{user_id}] = {v_i:.3f}（单位信息成本）

## 市场环境

**用户总数**：n = {n}

**信息相关系数**：ρ = {rho:.2f}
- 你的信息与其他用户相关
- ρ = 0：他人信息无法推断你
- ρ = 1：他人信息完美推断你
- 当前 ρ = {rho:.2f}，推断能力{rho_level}

**观测噪声**：σ² = {sigma_noise_sq}
- σ² 越大：噪声越大，泄露越低
- σ² 越小：数据越准确，泄露越高

**隐私偏好分布**：v ∈ [{v_min}, {v_max}]
- 你的位置：v = {v_i:.3f}（{v_description}）

## 核心机制

### 推断外部性

**关键洞察**：泄露信息量不仅取决于你是否分享，还取决于其他人是否分享。任何人的分享都会增加所有人（包括不分享者）的信息泄露量。

**如果你分享**：
- 获得补偿 p = {price:.4f}
- 你的信息从间接部分泄露 → 完全泄露

**如果你不分享**：
- 无补偿
- 保护未间接泄露的部分
- 但仍有基础泄露（他人分享导致）

### 次模性

**分享的人越多 → 你的边际泄露越小**
- 基础泄露越高 → 边际泄露越低
- 其他人分享得多 → 你再分享的额外成本减少

### 补偿逻辑

**平台报价旨在覆盖边际隐私损失**
- 报价 p = {price:.4f} 反映你的边际价值
- 你需要判断：p 是否足以覆盖 v × 边际泄露

## 决策框架

**隐私成本** = v × 边际泄露量

**权衡**：补偿收益 p vs 隐私成本 v × 边际泄露

## 输出

请输出严格JSON：
{{
  "share": 0或1（0=不分享，1=分享），
  "reason": "简要说明决策理由（不超过150字）"
}}
\end{verbatim}
  \end{minipage}}
\end{table*}

\begin{table*}[t]
  \centering
  \caption{Prompt template v5 (source: \texttt{b.v6}).}
  \label{tab:app_prompt_v5}
  \fbox{\begin{minipage}{0.97\textwidth}
\begin{verbatim}
# 场景：数据市场静态博弈（推断外部性）

你是用户 {user_id}，正在参与一个**一次性的数据市场决策**。

## 基本信息

**你的私有信息**：
- 你的隐私偏好：v[{user_id}] = {v_i:.3f}
- 平台给你的报价：p[{user_id}] = {price:.4f}

**公共知识**（所有人都知道）：
- 用户总数：n = {n}
- 类型相关系数：ρ = {rho:.2f}
  （你的类型与其他用户的类型相关，相关系数为 {rho:.2f}）
- 观测噪声：σ² = {sigma_noise_sq}
- 隐私偏好分布：所有用户的 v 均匀分布在 [{v_min}, {v_max}]
  （你的 v = {v_i:.3f}，相对位置：{v_description}，属于{v_level}隐私偏好群体）

**你不知道的信息**：
- 其他用户的具体 v 值（你只知道分布）
- 其他用户会如何决策（因为是同时决策）

## 推断外部性机制

**关键概念**：即使你不分享数据，平台也能通过其他人的数据推断你的信息。

**泄露信息量 I_i(a)**：
- 给定分享集合 S，平台对你的推断精度提升量
- 通过贝叶斯更新计算：I_i(S) = σ_i² - Var(X_i | S)
- **核心外部性**：I_i 不仅取决于你是否分享，还取决于其他人是否分享

**你的效用函数**：
- 如果你**分享**：u_i = p_i - v_i × I_i(你分享, 其他人的决策)
- 如果你**不分享**：u_i = 0 - v_i × I_i(你不分享, 其他人的决策)

**关键洞察**：
- 不分享也会有**基础泄露**（因为其他人分享会泄露你的信息）
- 分享的真正成本是**边际泄露** = I_i(分享) - I_i(不分享)
- 补偿价格 p_i 旨在覆盖你的边际隐私损失

## 理性预期决策框架

因为你不知道其他人会如何选择，你需要：

**1. 基于分布推测其他人的行为**：
- v 值较低的用户更可能分享（隐私成本低）
- v 值较高的用户更不可能分享（隐私成本高）
- 你的 v = {v_i:.3f}，处于{v_level}水平

**2. 计算期望效用**：
- E[u_i | 分享] = E[p_i - v_i × I_i(1, a_{{-i}})]
- E[u_i | 不分享] = E[- v_i × I_i(0, a_{{-i}})]

**3. 理解次模性**：
- 分享的人越多，你的边际信息价值越低
- 基础泄露越高（别人分享多），你分享的边际泄露越小

**4. 做出最佳反应**：
- 如果 E[u_i | 分享] > E[u_i | 不分享]，则分享
- 否则不分享

## 你的任务

基于上述机制，在**不知道其他人具体决策**的情况下，通过**理性预期**判断是否分享数据。

**思考要点**：
1. 你的 v 值在分布中的位置如何？（v = {v_i:.3f}，属于{v_level}群体）
2. 预期会有多少比例的用户分享？
3. 在那个预期下，你分享的边际价值是多少？
4. 报价 p = {price:.4f} 能否覆盖你的边际隐私损失？
5. 相关性 ρ = {rho:.2f} 如何影响外部性？

## 输出格式

请输出严格JSON：
{{
  "share": 0或1（0=不分享，1=分享），
  "reason": "简要说明你的权衡与信念依据（不超过150字）"
}}
\end{verbatim}
  \end{minipage}}
\end{table*}
